\documentclass[12pt]{article}
\usepackage{amsmath, amssymb, amsthm}
\usepackage[margin=1in]{geometry}
\usepackage{enumitem}
\usepackage{hyperref}

\title{Lattice Propagator and Continuum Limit for a Scalar Field on a Hypercubic Lattice}
\author{Student Code: 403200498}
\date{\today}

\begin{document}

\maketitle

\section*{Pipeline of the proof}
\begin{enumerate}[label=\textbf{\arabic*.}]
    \item Define the lattice and the lattice propagator.
    \item Derive the recursion relation for the propagator.
    \item Introduce the Brillouin zone and Fourier transform.
    \item Solve the recursion in momentum space to obtain the exact lattice propagator.
    \item Analyze convergence at fixed bare mass \(m_0\).
    \item Extract the correlation length and derive the tuning condition.
    \item Compute the required scaling of the bare mass \(m_0(\Delta)\).
    \item Match to the continuum Euclidean propagator and identify the physical mass.
\end{enumerate}

\section{Step 1: Define the lattice and the lattice propagator}

We work on a \(4\)-dimensional hypercubic lattice with lattice spacing \(\Delta\).  
Lattice points are labelled by integer 4-vectors
\[
n=(n_1,n_2,n_3,n_4)\in\mathbb Z^4 .
\]
The physical coordinate of the lattice point \(n\) is
\[
x=\Delta n .
\]
A \textbf{link} is a step from a lattice point to one of its \(8\) nearest neighbours:
\[
\pm e_\mu,\qquad \mu=1,2,3,4,
\]
where \(e_\mu\) is the unit vector in the \(\mu\)-th direction.  
A \textbf{path} from \(0\) to \(n\) is any sequence of such links whose net displacement is \(n\).  
The path may wander, go backward, cross itself, and may traverse a given link any number of times.

Let
\[
L(\text{path})=\text{number of links in the path}.
\]
Each link contributes a factor
\[
t=e^{-m_0},
\]
where \(m_0\) is the bare mass parameter in lattice units.  
Therefore a path of length \(L\) contributes
\[
t^L=e^{-m_0 L}.
\]
The \textbf{lattice propagator} from \(0\) to \(n\) is defined as
\begin{equation}
\boxed{
G(n)=\sum_{\text{lattice paths }0\to n} t^{L(\text{path})}.
}
\tag{1}
\end{equation}
Equivalently, since \(x=\Delta n\),
\[
G(n)\equiv G(x/\Delta).
\]

\section{Step 2: Derive the recursion relation for the propagator}

Split all paths from \(0\) to \(n\) into two classes.

\subsection{The empty path}
The empty path exists only when \(n=0\).  
It has length \(0\) and contributes
\[
1=\delta_{n,0}.
\]

\subsection{Non-empty paths}
Every non-empty path has a first step.  
The first step must be one of
\[
s e_\mu,\qquad s=\pm 1,\quad \mu=1,2,3,4.
\]
Suppose the first step is \(s e_\mu\).  
After this step, the path is at the lattice point \(s e_\mu\).  
The remaining part of the path must go from \(s e_\mu\) to \(n\).

By \textbf{translation invariance of the lattice}, the number and weights of the remaining paths from \(s e_\mu\) to \(n\) are the same as those of paths from \(0\) to \(n-s e_\mu\).  
Therefore the total contribution of all paths whose first step is \(s e_\mu\) is
\[
t\,G(n-s e_\mu).
\]
Summing over all possible first steps \(s e_\mu\) gives
\[
t\sum_{\mu=1}^4\bigl[G(n-e_\mu)+G(n+e_\mu)\bigr].
\]
Adding the empty-path contribution, we obtain the recursion relation
\begin{equation}
\boxed{
G(n)=\delta_{n,0}
+
t\sum_{\mu=1}^4\bigl[G(n-e_\mu)+G(n+e_\mu)\bigr].
}
\tag{2}
\end{equation}

\section{Step 3: Introduce the Brillouin zone and Fourier transform}

The \textbf{Brillouin zone} is the fundamental domain of momentum space for a system defined on a lattice. It is the set of all distinct momenta \(k\) that label the independent Fourier modes on the lattice.

Because the lattice is discrete, momenta are only defined modulo \(2\pi/\Delta\) (in units where \(\Delta=1\), modulo \(2\pi\)). This means that two momenta differing by \(2\pi\) in any component describe the same physical mode. Therefore, we can restrict \(k_\mu\) to an interval of length \(2\pi\), conventionally chosen as
\[
k_\mu\in[-\pi,\pi].
\]
This interval is the Brillouin zone for a hypercubic lattice with unit lattice spacing.

We use the representation of the Kronecker delta:
\begin{equation}
\delta_{n,0}
=
\int_{-\pi}^{\pi}\frac{d^4k}{(2\pi)^4}\,e^{ik\cdot n},
\qquad
k\cdot n=\sum_{\mu=1}^4 k_\mu n_\mu.
\tag{3}
\end{equation}
Write the propagator as a Fourier integral:
\begin{equation}
G(n)=
\int_{-\pi}^{\pi}\frac{d^4k}{(2\pi)^4}\,
e^{ik\cdot n}\,\widetilde G(k).
\tag{4}
\end{equation}
Substitute (3) and (4) into the recursion (2):
\[
\int\frac{d^4k}{(2\pi)^4}\,e^{ik\cdot n}\,\widetilde G(k)
=
\int\frac{d^4k}{(2\pi)^4}\,e^{ik\cdot n}
\left[
1
+
t\sum_{\mu=1}^4
\left(
e^{-ik_\mu}+e^{ik_\mu}
\right)
\widetilde G(k)
\right].
\]
Since this must hold for all \(n\), the integrands must be equal:
\[
\widetilde G(k)
=
1
+
2t\sum_{\mu=1}^4\cos k_\mu\;\widetilde G(k).
\]
Solving for \(\widetilde G(k)\),
\begin{equation}
\widetilde G(k)
=
\frac{1}{
1-2t\sum_{\mu=1}^4\cos k_\mu
}.
\tag{5}
\end{equation}

\section{Step 4: Exact lattice propagator}

Therefore the exact lattice propagator is
\begin{equation}
\boxed{
G(n)
=
\int_{-\pi}^{\pi}\frac{d^4k}{(2\pi)^4}\,
\frac{e^{ik\cdot n}}{
1-2t\sum_{\mu=1}^4\cos k_\mu
}.
}
\tag{6}
\end{equation}
Define the denominator
\begin{equation}
D(k)=1-2t\sum_{\mu=1}^4\cos k_\mu.
\tag{7}
\end{equation}
Then
\begin{equation}
G(n)=
\int_{-\pi}^{\pi}\frac{d^4k}{(2\pi)^4}\,
\frac{e^{ik\cdot n}}{D(k)}.
\tag{8}
\end{equation}

\section{Step 5: Convergence at fixed \(m_0\)}

At \(k=0\),
\begin{equation}
D(0)=1-8t.
\tag{9}
\end{equation}

\subsection{Case \(m_0<\log 8\)}
Then
\[
t=e^{-m_0}>\frac18,
\qquad
D(0)=1-8t<0.
\]
As \(k\) varies over the Brillouin zone, \(\sum_{\mu=1}^4\cos k_\mu\) ranges from \(-4\) to \(4\).  
Thus \(D(k)\) ranges from \(1-8t<0\) to \(1+8t>0\).  
By continuity, \(D(k)\) vanishes somewhere in the Brillouin zone.  
Therefore the integrand in (8) has a pole on the integration contour, and the sum over paths diverges.

\subsection{Case \(m_0>\log 8\)}
Then
\[
t<\frac18,
\qquad
D(k)>0\quad\text{for all }k.
\]
The integral converges.  
For large \(|n|\), the propagator decays exponentially with a finite correlation length in lattice units.

\subsection{Case \(m_0=\log 8\)}
Then
\[
t=\frac18,
\qquad
D(k)\ge 0,
\]
with \(D(k)=0\) only at \(k=0\).  
In \(4\) dimensions the integral converges, but \(G(n)\) decays only as a power.

\subsection{Conclusion for fixed \(m_0\)}
At fixed \(m_0\), in the limit \(\Delta\to 0\) with \(x\) fixed:
\begin{itemize}
    \item If \(m_0<\log 8\), the sum diverges.
    \item If \(m_0\ge \log 8\), the sum converges but is dominated by paths of length proportional to \(|n|\), and \(G(n)\to 0\) for fixed physical separation \(x\).
    \item To keep long paths relevant in the continuum limit, \(m_0\) must be tuned.
\end{itemize}

\section{Step 6: Correlation length and tuning condition}

For \(t<1/8\), \(D(k)>0\) on the real contour.  
We may deform the \(k\)-contour into the complex plane by Cauchy's theorem, because the integrand is analytic except at poles.

To extract the large-\(|n|\) behaviour, write
\[
n=|n|\,\hat n,
\qquad
\hat n_\mu=\frac{n_\mu}{|n|},
\qquad
\sum_{\mu=1}^4\hat n_\mu^2=1.
\]
Set
\[
k=i\kappa\,\hat n,
\qquad \kappa>0.
\]
Then
\[
e^{ik\cdot n}=e^{i(i\kappa\hat n)\cdot n}
=e^{-\kappa|n|}.
\]
The leading exponential decay comes from the pole of \(1/D(k)\) closest to the real axis.  
The pole condition is
\[
D(i\kappa\hat n)=0.
\]
Using (7),
\begin{equation}
D(i\kappa\hat n)
=
1-2t\sum_{\mu=1}^4\cosh(\kappa\hat n_\mu).
\tag{10}
\end{equation}
For small \(\kappa\), expand
\[
\cosh(\kappa\hat n_\mu)
=
1+\frac12\kappa^2\hat n_\mu^2+O(\kappa^4).
\]
Therefore
\[
\sum_{\mu=1}^4\cosh(\kappa\hat n_\mu)
=
4+\frac12\kappa^2\sum_{\mu=1}^4\hat n_\mu^2+O(\kappa^4)
=
4+\frac12\kappa^2+O(\kappa^4).
\]
Substitute into (10):
\begin{equation}
D(i\kappa\hat n)
=
1-2t\left(4+\frac12\kappa^2+O(\kappa^4)\right)
=
1-8t-t\kappa^2+O(\kappa^4).
\tag{11}
\end{equation}
Setting this to zero gives
\[
1-8t-t\kappa^2=0,
\]
so
\begin{equation}
\kappa^2=\frac{1-8t}{t}.
\tag{12}
\end{equation}
Therefore the correlation length in lattice units is
\begin{equation}
\xi_{\rm lat}=\frac{1}{\kappa}
=
\sqrt{\frac{t}{1-8t}}.
\tag{13}
\end{equation}
For \(t\) close to \(1/8\), write \(t=1/8-\varepsilon\). Then \(1-8t=8\varepsilon\), and
\[
\xi_{\rm lat}
=
\sqrt{\frac{1/8-\varepsilon}{8\varepsilon}}
\approx
\frac{1}{8\sqrt{\varepsilon}}
=
\frac{1}{8\sqrt{(1-8t)/8}}
=
\frac{\sqrt{2}}{4}\frac{1}{\sqrt{1-8t}}.
\]
Thus, up to a constant factor,
\begin{equation}
\boxed{
\xi_{\rm lat}\sim \frac{1}{\sqrt{1-8t}}.
}
\tag{14}
\end{equation}
The physical correlation length is
\begin{equation}
\xi_{\rm phys}=\Delta\,\xi_{\rm lat}
\sim
\frac{\Delta}{\sqrt{1-8t}}.
\tag{15}
\end{equation}
To keep \(\xi_{\rm phys}\) finite as \(\Delta\to 0\), we need
\[
\sqrt{1-8t}\sim \Delta.
\]
Squaring,
\[
1-8t\sim \Delta^2.
\]
Therefore we write
\begin{equation}
\boxed{
1-8t=8c\,\Delta^2+O(\Delta^4).
}
\tag{16}
\end{equation}
Solving for \(t\),
\begin{equation}
\boxed{
t=\frac18-c\,\Delta^2+O(\Delta^4).
}
\tag{17}
\end{equation}

\section{Step 7: Compute the required scaling of the bare mass \(m_0(\Delta)\)}

We have
\[
m_0=-\log t.
\]
Substitute (17):
\[
t=\frac18\left(1-8c\,\Delta^2+O(\Delta^4)\right).
\]
Then
\[
m_0
=
-\log\left[\frac18\left(1-8c\,\Delta^2+O(\Delta^4)\right)\right]
=
\log 8
-
\log\left(1-8c\,\Delta^2+O(\Delta^4)\right).
\]
Use the expansion
\[
-\log(1-x)=x+\frac{x^2}{2}+O(x^3),
\qquad x\to 0.
\]
Here \(x=8c\,\Delta^2+O(\Delta^4)\). Therefore
\begin{equation}
m_0
=
\log 8
+
8c\,\Delta^2
+O(\Delta^4).
\tag{18}
\end{equation}
So
\begin{equation}
\boxed{
m_0(\Delta)
=
\log 8+8c\,\Delta^2+O(\Delta^4).
}
\tag{19}
\end{equation}
Equivalently,
\begin{equation}
\boxed{
m_0(\Delta)=\log 8+O(\Delta^2).
}
\tag{20}
\end{equation}

\section{Step 8: Match to the continuum Euclidean propagator and identify the physical mass}

Expand \(D(k)\) for small \(k\) and small \(\Delta\):
\[
D(k)=1-2t\sum_{\mu=1}^4\cos k_\mu.
\]
For small \(k\),
\[
\sum_{\mu=1}^4\cos k_\mu
=
4-\frac12 k^2+O(k^4),
\qquad
k^2=\sum_{\mu=1}^4 k_\mu^2.
\]
Using (17),
\[
2t\sum_{\mu=1}^4\cos k_\mu
=
2\left(\frac18-c\,\Delta^2\right)
\left(4-\frac12 k^2\right)
+O(\Delta^4,k^4)
=
1-\frac18 k^2-8c\,\Delta^2+O(\Delta^4,k^4).
\]
Therefore
\begin{equation}
D(k)
=
1-\left(1-\frac18 k^2-8c\,\Delta^2+O(\Delta^4,k^4)\right)
=
\frac18 k^2+8c\,\Delta^2+O(\Delta^4,k^4).
\tag{21}
\end{equation}
Rescale the momentum:
\[
k=\Delta p.
\]
Then
\begin{equation}
D(k)
=
\Delta^2\left(\frac18 p^2+8c\right)+O(\Delta^4).
\tag{22}
\end{equation}
The lattice propagator becomes
\[
G(n)
=
\int_{-\pi}^{\pi}\frac{d^4k}{(2\pi)^4}\,
\frac{e^{ik\cdot n}}{D(k)}
\approx
\int
\frac{\Delta^4 d^4p}{(2\pi)^4}\,
\frac{e^{ip\cdot x}}{
\Delta^2\left(\frac18 p^2+8c\right)
}.
\]
Simplifying,
\begin{equation}
G(n)
\approx
8\Delta^2
\int
\frac{d^4p}{(2\pi)^4}\,
\frac{e^{ip\cdot x}}{
p^2+64c
}.
\tag{23}
\end{equation}
The continuum Euclidean propagator for a scalar field of mass \(m\) is
\begin{equation}
D_E(x)
=
\int
\frac{d^4p}{(2\pi)^4}\,
\frac{e^{ip\cdot x}}{
p^2+m^2
}.
\tag{24}
\end{equation}
Comparing (23) and (24), we identify
\begin{equation}
\boxed{
m^2=64c.
}
\tag{25}
\end{equation}
Thus
\[
c=\frac{m^2}{64},
\qquad
8c=\frac{m^2}{8}.
\]
Substituting into (19),
\begin{equation}
\boxed{
m_0(\Delta)
=
\log 8+\frac{m^2}{8}\Delta^2+O(\Delta^4).
}
\tag{26}
\end{equation}
The overall factor \(8\Delta^2\) in (23) is absorbed into the functional measure of the continuum path integral.

\section*{Final answer}

The lattice propagator is
\[
G(n)=\sum_{\text{paths }0\to n} t^{L(\text{path})},
\qquad
t=e^{-m_0}.
\]
It satisfies the recursion
\[
G(n)=\delta_{n,0}
+
t\sum_{\mu=1}^4\bigl[G(n-e_\mu)+G(n+e_\mu)\bigr].
\]
Its exact Fourier representation is
\[
G(n)=
\int_{-\pi}^{\pi}\frac{d^4k}{(2\pi)^4}\,
\frac{e^{ik\cdot n}}{
1-2t\sum_{\mu=1}^4\cos k_\mu
}.
\]
At fixed \(m_0\), as \(\Delta\to 0\):
\begin{itemize}
    \item if \(m_0<\log 8\), the sum diverges;
    \item if \(m_0\ge \log 8\), the sum is dominated by paths of length proportional to \(|n|\), and the propagator vanishes for fixed physical separation \(x\).
\end{itemize}
To keep long paths relevant in the continuum limit, the bare mass must be tuned as
\[
\boxed{
m_0(\Delta)=\log 8+O(\Delta^2).
}
\]
More precisely,
\[
\boxed{
m_0(\Delta)
=
\log 8+\frac{m^2}{8}\Delta^2+O(\Delta^4),
}
\]
where \(m\) is the physical mass. With this tuning, the lattice propagator becomes proportional to the continuum Euclidean propagator
\[
D_E(x)=\int\frac{d^4p}{(2\pi)^4}\frac{e^{ipx}}{p^2+m^2},
\]
up to an overall normalization factor absorbed into the functional measure.

\subsection*{AI Use Declaration}

\textbf{Model(s) used:} DeepSeek.

\medskip
\noindent\textbf{Revisions and corrections made to model output, and clarifications added:}
\begin{enumerate}[label=\arabic*.]
    \item In Step~1, the hypercubic lattice is now defined with full detail, including the physical coordinate \(x=\Delta n\). The first model output did not define the lattice details clearly.
    \item In the model output, the recursion formula was introduced suddenly without explaining its origin. It is now fully justified in Step~2.
    \item The model used the momentum interval \(k_\mu\in[-\pi,\pi]\) without explanation or justification. This is now justified in Step~3.
    \item In the first version, convergence was not checked and its region was not determined. This is now explained in detail in Step~5.
    \item In the first version, the model obtained the general form of \(m_0\) by jumping over several relations and details. In this version, it is derived in Steps~6 and~7 without skipping any details.
    \item In the first version, the model did not show why the constant \(c\) in Eq.~(19) is chosen as \(m^2/64\). This is now shown explicitly in Step~8 by matching the general form with the Euclidean propagator.
\end{enumerate}

\medskip
\noindent\textbf{Comparisons and checks against the lecture notes:}
\begin{enumerate}[label=\arabic*.]
    \item Choice of lattice paths: in the last paragraph of p.~56 and the first paragraph of p.~57, the lecture proposes the hypercubic lattice paths.
    \item Definition of the recursion formula: the lecture defines a recursion formula on p.~58, Eq.~(7.3).This is a insight that define a recursion formula can help to derive the propagator 
\end{enumerate}

\medskip
\noindent\textbf{Calculation steps, limiting cases, and consistency checks verified by hand:}
I did not use hand calculations for checking. I checked the logic and correctness of the reasoning by eye and by mental verification. The correctness of the steps and the coherence between steps were checked by mind and AI. I tried my best to avoid jumping between steps and relations without proving and justifying them.

\medskip
\noindent\textbf{Calculation steps, limiting cases, and consistency checks subjected to automated verification (via Mathematica, Python, LLMs):}
I explained all edits in the ``Revisions and corrections'' part above. The consistency checks were done by eye and mind.

\end{document}
